\section{Escape estimates and compact superlevels}
\label{sec:escape}

The geometric argument will compare critical points above a fixed positive
level.  We first show that such points cannot escape to the boundary or to
infinity for almost every target value.  The main step near the real axis is
a bound on the total length of a horizontal line that maps into a compact
target set.

Throughout this section assume $\rho>2$, let $\beta=\log_2\rho>1$, and
let $P$ be the probability from Lemma~\ref{lem:weighted-law}.  Thus, for
$z=x+iy\in\HH$ and every nonnegative Borel function $\Phi$ on $\KK$,
\begin{equation}
 \EE\!\left[|Q_F(z)|^2\Phi(T_zF)\right]
       =y^{-\beta}\EE\Phi(F).
 \label{escape:covariance}
\end{equation}
Set
\begin{equation}
 k=\frac{\beta+3}{4}>1,\qquad
 V_\xi^F(x+iy)=\frac{y^k}{|F(x+iy)-\xi|}
 \quad(F\in\KK,\ \xi\in\CC),
 \label{escape:potential}
\end{equation}
with value $+\infty$ at the unique possible point where $F(z)=\xi$.

\begin{proposition}[Compact positive superlevels]
\label{prop:proper-superlevels}
For $P$-almost every $F\in\KK$, for planar area-almost every $\xi\in\CC$,
the set
\[
 \{z\in\HH:V_\xi^F(z)\ge\delta\}
\]
is compact in $\HH$ for every $\delta>0$.
\end{proposition}

We first obtain a weighted area integral from
\eqref{escape:covariance}.  Its tail gives escape at positive heights.
An integration by parts then controls horizontal preimage length, which
will bound the area of exceptional targets near the real axis.

\subsection{A weighted area integral}

Recall $m(z)=(z-i)/(z+i)$ and the affine inverse
$z^\#=(-x+i)/y$ of $z=x+iy$.  Define
\begin{equation}
 S_0=\{z\in\HH:|m(z)|\ge\tfrac12\},\qquad
 I(F)=\int_{S_0}y^{\beta+1}|F(x+iy)|^{-4}\dd A(x+iy).
 \label{escape:weighted-integral}
\end{equation}
The complement of $S_0$ has compact closure in $\HH$ and contains the
unique zero $i$ of every $F\in\KK$.

\begin{lemma}[Weighted inverse area]
\label{lem:weighted-inverse-area}
The function $I:\KK\to[0,\infty]$ is Borel and satisfies
$\EE I(F)\le256\pi$.  In particular, $I(F)<\infty$ almost surely.
\end{lemma}
\begin{proof}
For $z\ne i$, normalization gives
\[
 (T_zF)(z^\#)=-\frac{F(z)Q_F(z)}y.
\]
Apply \eqref{escape:covariance} with
$\Phi(H)=|H(z^\#)|^{-4}$.  Both evaluations avoid the unique zero, and
we obtain
\begin{equation}
 \EE\!\left[|F'(z)|^2|F(z)|^{-4}\right]
       =y^{-\beta-4}\EE|F(z^\#)|^{-4}.
 \label{escape:inverse-identity}
\end{equation}
The transformation $(x,y)\mapsto(u,Y)=(-x/y,1/y)$ is an involution, with
\[
 \det\frac{\partial(u,Y)}{\partial(x,y)}
 =\det\begin{pmatrix}-1/y&x/y^2\\0&-1/y^2\end{pmatrix}
 =y^{-3}.
\]
Consequently $dx\,dy=Y^{-3}du\,dY$.  Moreover,
\[
 |m(z^\#)|^2
 =\frac{x^2+(y-1)^2}{x^2+(y+1)^2}=|m(z)|^2,
\]
so the involution preserves $S_0$.

Lemma~\ref{lem:compact-class} gives $|F|\ge1/16$ on $S_0$.
Injectivity and conformal change of variables therefore imply
\begin{equation}
 \int_{S_0}|F'(z)|^2|F(z)|^{-4}\dd A(z)
 =\int_{F(S_0)}|w|^{-4}\dd A(w)
 \le\int_{|w|\ge1/16}|w|^{-4}\dd A(w)=256\pi.
 \label{escape:deterministic-area}
\end{equation}
All the integrands are nonnegative and jointly Borel, by continuity of
interior evaluations.  Integrating \eqref{escape:inverse-identity},
using Tonelli and the involution, gives
\[
 \begin{split}
 \EE\int_{S_0}|F'|^2|F|^{-4}\dd A
 &=\int_{S_0}y^{-\beta-4}\EE|F(z^\#)|^{-4}\,dx\,dy\\
 &=\int_{S_0}Y^{\beta+1}\EE|F(u+iY)|^{-4}\,du\,dY
 =\EE I(F).
 \end{split}
\]
The same joint measurability makes $I$ Borel.  The conclusion follows
from \eqref{escape:deterministic-area}.
\end{proof}

\subsection{Escape and horizontal preimage length}

\begin{lemma}[Escape and trace bounds]
\label{lem:escape-trace}
Let $F\in\KK$ satisfy $I(F)<\infty$.  Then
\begin{equation}
 \frac{y^k}{|F(x+iy)|}\longrightarrow0
       \qquad\text{as }|x|+y\longrightarrow\infty.
 \label{escape:tail}
\end{equation}
For every compact $K\subset\CC$ and $\epsilon>0$,
$F^{-1}(K)\cap\{\Im z\ge\epsilon\}$ is compact in $\HH$, and
\[
 H_{F,K}:=\sup\{\Im z:F(z)\in K\}<\infty,
\]
where the supremum of an empty set is assigned zero.  Moreover, there is
$C_{F,K}<\infty$ such that
\begin{equation}
 \int_{\mathbb R}\mathbf1_K(F(x+i\epsilon))\,dx
       \le C_{F,K}\qquad(\epsilon>0).
 \label{escape:trace}
\end{equation}
\end{lemma}
\begin{proof}
Write $R(z)=|x|+y$.  If $z'=x'+iy'\in B(z,y/2)$, then
\[
 y/2<y'<3y/2,\qquad
 R(z)\le |x'|+3y'\le3R(z').
\]
Thus these disks lie in $\{R(z')\ge R(z)/3\}$ and eventually in $S_0$.
On them $1/F$ is holomorphic.  Submean for $|1/F|^4$, followed by
comparison of $y$ with $y'$, gives
\begin{equation}
 y^{\beta+3}|F(z)|^{-4}
 \le C_\beta\!\int_{S_0\cap\{R(z')\ge R(z)/3\}}
       (\Im z')^{\beta+1}|F(z')|^{-4}\dd A(z')
 \longrightarrow0.
 \label{escape:uniform-tail}
\end{equation}
The convergence follows from the tail of the single integrable function
in \eqref{escape:weighted-integral}, and is uniform in $z$ as
$R(z)\to\infty$.  Since $4k=\beta+3$, this proves
\eqref{escape:tail}.

At heights $y\ge\epsilon>0$, \eqref{escape:tail} also gives
$|F(z)|\to\infty$ as $R(z)\to\infty$.  Hence the preimage of a compact
$K$, restricted to these heights, is bounded.  It is closed in the
plane, because a limit stays at height at least $\epsilon$ and
$F$ is continuous there.  It is therefore compact in $\HH$.
Taking $\epsilon=1$ and then including the points below height one
proves the bound on $H_{F,K}$.

To prove \eqref{escape:trace}, choose a nonnegative
$\psi\in C_c^\infty(\CC)$ with $\psi\ge1$ on $K$, and fix $\epsilon>0$.
The function $g=\psi\circ F$ has compact support when restricted to
$y\ge\epsilon/2$, by the compactness just proved.  Thus
\[
 h(y)=\int_{\mathbb R}g(x,y)\,dx
\]
is smooth for $y\ge\epsilon$, differentiation passes under the integral,
and $h,h'$ vanish at sufficiently large height.  The horizontal second
derivative of $g$ integrates to zero.  Therefore
\[
 \begin{split}
 \int_{y>\epsilon}(y-\epsilon)\Delta g(x,y)\,dx\,dy
 &=\int_\epsilon^\infty(y-\epsilon)h''(y)\,dy\\
 &=[(y-\epsilon)h'(y)]_\epsilon^\infty
       -[h(y)]_\epsilon^\infty=h(\epsilon).
 \end{split}
\]
Write $\Omega=F(\HH)$ and $Y(w)=\Im F^{-1}(w)$ for $w\in\Omega$.
Since $\Delta(\psi\circ F)=|F'|^2(\Delta\psi)\circ F$,
conformal change of variables gives
\[
 h(\epsilon)=
 \int_{\Omega\cap\{Y>\epsilon\}}(Y(w)-\epsilon)
                    \Delta\psi(w)\dd A(w).
\]
Taking the absolute value of this signed integral yields
\[
 \int_{\mathbb R}\mathbf1_K(F(x+i\epsilon))\,dx
 \le h(\epsilon)
 \le H_{F,\operatorname{supp}\psi}
                       \|\Delta\psi\|_{L^1(\CC)}.
\]
This finite bound is independent of $\epsilon$, as required.
\end{proof}

All conclusions of Lemma~\ref{lem:escape-trace} hold on the single
full-measure set $\{I(F)<\infty\}$, simultaneously for every compact
target set.  In particular, no exceptional set depending on the height
is needed.  We now use the trace bound to control targets approached at
small height.

\subsection{Targets near the real axis}

\begin{proof}[Proof of Proposition~\ref{prop:proper-superlevels}]
Fix $F$ with $I(F)<\infty$, and put $\Omega=F(\HH)$.
For each fixed $\xi\in\CC$ and $\epsilon>0$, Lemma~\ref{lem:escape-trace}
implies
\begin{equation}
 V_\xi^F(x+iy)\longrightarrow0
 \quad\text{as }|x|+y\longrightarrow\infty,\quad y\ge\epsilon.
 \label{escape:target-tail}
\end{equation}
Indeed $|F|\to\infty$ in this region, so eventually
$|F-\xi|\ge|F|/2$.

For small heights first consider omitted targets.  Fix $R,\delta>0$ and
a dyadic $h\le1$.  Partition $\mathbb R$ into the half-open cells
$I_j=[jh,(j+1)h)$, $j\in\mathbb Z$.  Call a cell a witness cell if there
exist $x\in I_j$, $h\le y\le2h$, and
$\xi\notin\Omega$ with $|\xi|\le R$ such that
$V_\xi^F(x+iy)\ge\delta$.
At any such witness $z=x+iy$,
\[
 |F(z)-\xi|\le(2h)^k/\delta.
\]
The disk $B(z,y)$ lies in $\HH$, and its image omits $\xi$.
Lemma~\ref{lem:classical-univalence} therefore gives
\begin{equation}
 y|F'(z)|\le4|F(z)-\xi|\le4(2h)^k/\delta.
 \label{escape:witness-derivative}
\end{equation}
For every $x'\in I_j$, the relative point
\[
 (x'-x)/y+i h/y
       \in[-1,1]+i[1/2,1]
\]
belongs to a fixed compact subset of $\HH$.
The values of $T_zF\in\KK$ there are bounded by a common constant $M$,
by Lemma~\ref{lem:compact-class}.  The definition of $T_zF$ and
\eqref{escape:witness-derivative} imply
\begin{equation}
 |F(x'+ih)-\xi|
 \le (1+4M)(2h)^k/\delta
 =:C_\delta h^k.
 \label{escape:witness-cell}
\end{equation}

In particular, each witness cell at height $h$ maps entirely into the
fixed target disk $K_{R,\delta}=\{|w|\le R+C_\delta\}$.
By \eqref{escape:trace}, the disjoint cells, each of length $h$, have
number at most $C_{F,R,\delta}/h$.  This count permits all the omitted
targets with $|\xi|\le R$ at once.
For every witness cell place a closed target disk of radius $C_\delta h^k$
at the image of its midpoint at height $h$.
Estimate~\eqref{escape:witness-cell}, applied to any witness in that
cell, shows that these disks cover every target witnessed at scale $h$.
Their union $E_h$ has area at most
\begin{equation}
 |E_h|\le C_{F,R,\delta}h^{2k-1}.
 \label{escape:target-cover}
\end{equation}
The sum over dyadic $h\to0$ is finite.  Hence
$\limsup_{n\to\infty}E_{2^{-n}}$ has area zero: its area is bounded by
each tail of that convergent sum.

Take the countable union of these null sets for positive integer $R$
and $\delta=1/n$, $n\ge1$.  Every omitted target outside the union
satisfies
\begin{equation}
 \sup_{x\in\mathbb R}V_\xi^F(x+iy)\longrightarrow0
                   \qquad(y\downarrow0).
 \label{escape:small-height}
\end{equation}
Indeed each positive threshold eventually has no witness in any smaller
dyadic slab, and these slabs cover all sufficiently small heights.
This argument only uses finite disk covers for the fixed map $F$.

For an interior target $\xi\in\Omega$, choose $a>0$ so that
$\overline B(\xi,a)\subset\Omega$.  Its image under the continuous
inverse $F^{-1}$ is compact in $\HH$, and has positive minimum height.
Below that height, $|F(z)-\xi|\ge a$, so
$V_\xi^F(z)\le y^k/a$.  Thus \eqref{escape:small-height} holds for every
interior target as well.

For each target retained above and each $\delta>0$,
\eqref{escape:small-height} first excludes all sufficiently small
heights from $\{V_\xi^F\ge\delta\}$.
Equation~\eqref{escape:target-tail} then excludes all sufficiently
large $|x|+y$.  Finally
\[
 \{V_\xi^F\ge\delta\}
 =\{z\in\HH:|F(z)-\xi|\le(\Im z)^k/\delta\}
\]
is closed, includes the possible pole, and lies in a bounded box with
positive minimum height.  It is compact in $\HH$.  The assertion follows
from Lemma~\ref{lem:weighted-inverse-area}.
\end{proof}

\begin{lemma}[Measurability of compact superlevels]
\label{lem:superlevel-borel}
The set of pairs $(F,\xi)\in\KK\times\CC$ for which every closed positive
superlevel of $V_\xi^F$ is compact in $\HH$ is Borel.
\end{lemma}
\begin{proof}
Use the compact exhaustion
$C_m=\{x+iy:|x|\le m,\ 1/m\le y\le m\}$, $m\ge1$.
The stated property is equivalent to
\[
 \text{for every }n\ge1\text{ there is }m\ge1
       \text{ such that }\{V_\xi^F>1/n\}\subset C_m.
\]
Necessity follows by containing the closed $1/n$-superlevel in a box.
For sufficiency, given $\delta>0$ choose $1/n<\delta$; the corresponding
box contains the closed $\delta$-superlevel.
For fixed $n,m$, the containment is equivalent to
\[
 |F(q)-\xi|\ge n(\Im q)^k
 \quad\text{for every }q\in(\mathbb Q+i\mathbb Q)\cap(\HH\setminus C_m).
\]
Any violation occurs in a nonempty open subset of the open complement
of $C_m$ and is therefore detected at a rational point.
Each displayed inequality is closed in the parameters by continuity
of evaluation.  The countable tests prove the claim.
\end{proof}
